%
% teil1.tex -- Beispiel-File für das Paper
%
% (c) 2020 Prof Dr Andreas Müller, Hochschule Rapperswil
%
% !TEX root = ../../buch.tex
% !TEX encoding = UTF-8
%
\section{Eulers Produktdarstellung des Sinus}
\kopfrechts{Eulers Produktdarstellung des Sinus}

\subsection{Der Ansatz}

Euler beobachtete, dass die Nullstellen von $\sin(x)$ genau die Punkte
$x = n\pi$, $n \in \mathbb{Z}$, sind.
Analog zur Polynomfaktorisierung schrieb er einen formalen Ansatz ohne Konvergenzbedenken.
Sein Ansatz hatte die Form
\begin{equation}
\sin(x)
=
x \cdot \prod_{n=1}^{\infty} \biggl(1 - \frac{x^2}{n^2 \pi^2}\biggr).
\end{equation}
Jeder Faktor $(1 - x^2/n^2\pi^2)$ hat genau die Nullstellen $x = \pm n\pi$.
Der Vorfaktor $x$ erzeugt die Nullstelle bei $x = 0$.
Das Normierungsverhalten $\sin(x)/x \to 1$ für $x \to 0$ ist automatisch erfüllt,
da jeder Faktor für $x = 0$ den Wert $1$ annimmt.

Eulers Argument stützte sich auf die Analogie mit Polynomen und auf die
Übereinstimmung der Nullstellenstruktur.
Einen strengen Beweis lieferte er nicht, dieser wurde erst durch Weierstraß erbracht.

\subsection{Das Basel-Problem}

Dieser Abschnitt behandelt das Basel-Problem nur oberflächlich.
\index{Basel-Problem}%
Für eine detaillierte Analyse siehe Kapitel \ref{chapter:basel}.
Ein bemerkenswertes Nebenresultat ergibt sich durch Koeffizientenvergleich.
Die Taylorreihe des Sinus liefert den Koeffizienten von $x^3$ als $-1/6$.
Multipliziert man das unendliche Produkt formal aus, trägt jeder Faktor
$(1 - x^2/n^2\pi^2)$ den Term $-x^2/n^2\pi^2$ bei,
sodass der Gesamtkoeffizient von $x^3$ gleich $-\sum_n 1/(n^2\pi^2)$ ist.
Gleichsetzen liefert das berühmte Resultat:
\begin{equation}
\sum_{n=1}^{\infty} \frac{1}{n^2}
=
\frac{\pi^2}{6}.
\end{equation}
Dieses, als Basel-Problem bekannte Resultat, war seit Jahrzehnten offen
und hatte namhafte Mathematiker wie die Bernoullis beschäftigt.
\index{Bernoulli}%
Der Schritt ist analog zur endlichen Situation:
Multipliziert man $(1+a_1)(1+a_2)\cdots$ aus, so liefert jeder Faktor genau
einen linearen Beitrag $a_i$;
Produkte mehrerer solcher Terme tragen zu höheren Potenzen bei und spielen
beim Vergleich des $x^2$-Koeffizienten keine Rolle.
Eulers Lösung, beschrieben in \cite{produkt:euler}, war eine Sensation.
Stein und Shakarchi bemerken in \textit{Complex Analysis} \cite{produkt:steinShakarchi} treffend,
dass Eulers Herleitung zwar nicht rigoros war, aber den richtigen Weg wies,
der später durch die Weierstraß-Theorie vollständig gerechtfertigt werden sollte.

\section{Das Problem der Konvergenz}
\kopfrechts{Das Problem der Konvergenz}

Eulers Vorgehen war kühn, aber mathematisch nicht gesichert.
Zwei wesentliche Fragen bleiben offen.

\begin{enumerate}
\item \emph{Konvergiert das unendliche Produkt
$\prod(1 - x^2/n^2\pi^2)$ überhaupt?}
\index{Konvergenz}%
Im vorliegenden Fall lässt sich dies bejahen, da $\sum 1/n^2$ konvergiert.
Im Allgemeinen konvergiert ein unendliches Produkt $\prod(1+a_n)$ unter
geeigneten Bedingungen (etwa wenn $\sum |a_n|$ konvergiert) gegen einen Wert $\neq 0$;
die einfache Äquivalenz zur Konvergenz von $\sum a_n$ gilt nur in dieser
Einschränkung \cite{produkt:remmert}.

\item \emph{Folgt aus gleichen Nullstellen auch Gleichheit der Funktionen?}
Besitzen zwei ganze Funktionen dieselben Nullstellen (inklusive ihrer Vielfachheiten),
so ist ihr Quotient eine nullstellenfreie ganze Funktion.
Nach den Sätzen der Funktionentheorie lässt sich eine solche stets in der Form
$e^{g(z)}$ schreiben, wobei $g$ eine ganze Funktion ist.
Zwei ganze Funktionen mit denselben Nullstellen unterscheiden sich folglich
nur um diesen nullstellenfreien Faktor $e^{g(z)}$.
Euler nutzte implizit, dass dieser Faktor für den Sinus gleich $1$ ist,
was er durch das Verhalten für $x \to 0$ heuristisch erschloss,
aber nicht rigoros begründete.
\end{enumerate}

Diese Lücken machen deutlich, dass eine allgemeine Theorie benötigt wird:
eine Theorie, die präzise angibt, wann und wie eine ganze Funktion als Produkt
über ihre Nullstellen darstellbar ist.
Diese Theorie lieferte Weierstraß.
